\section{Introduction}
\label{sec:intro}

Modern data platforms increasingly expose Large Language Models (LLMs) as declarative semantic operators over unstructured text, spanning commercial systems such as BigQuery~\cite{google2026bigqueryai}, Snowflake Cortex AI~\cite{snowflake2026cortexai}, and PolarDB~\cite{alibaba2026polardbai}, alongside research engines~\cite{DBLP:journals/pvldb/PatelJPGAGZ25,DBLP:journals/corr/abs-2603-11622,DBLP:journals/pvldb/WangLWXL25}. Among these operators, semantic \textsf{ORDER BY} (with optional \LIMIT~$K$) orders records by natural-language criteria—such as review helpfulness, evidential support, or passage relevance in RAG~\cite{DBLP:conf/nips/YuP0WYZSC24}. Unlike row-independent semantic filters~\cite{DBLP:journals/pacmmod/ZeighamiSP25,DBLP:journals/corr/abs-2509-12610}, semantic \textsf{ORDER BY} couples all $n$ tuples globally under an expensive, context-bounded, and noisy LLM comparator. Each request repeats the instruction header alongside every compared record, so reducing call count alone can still inflate input tokens if long tuples are repeatedly resubmitted.

Selecting an LLM comparison primitive already forces a fundamental quality--efficiency trade-off. \emph{Pointwise} scoring~\cite{DBLP:conf/emnlp/DrozdovZD0RWAIM23,DBLP:conf/emnlp/SachanLJAYPZ22} minimizes requests ($O(n)$ or $O(n/m)$), but isolated prompts lack shared comparative context, relying on parametric world knowledge and degrading on relative semantic judgments. \emph{Pairwise}~\cite{DBLP:conf/acl/LuoCH024,DBLP:conf/naacl/QinJHZWYSLLMWB24,DBLP:journals/pvldb/PatelJPGAGZ25} and small-arity \emph{setwise}~\cite{DBLP:conf/sigir/ZhuangZKZ24} methods compare items directly, yet underutilize the model's effective window $m$ and pay fixed prompt headers across $O(n\log n)$ calls. Bounded \emph{listwise} prompting~\cite{DBLP:conf/emnlp/0001YMWRCYR23,DBLP:journals/corr/abs-2305-02156} resolves both limits by calibrating up to $m$ tuples jointly in one prompt and extracting $\binom{m}{2}$ pairwise preferences per call. For $n > m$, treating this capacity-$m$ ranker as a black-box primitive reduces physical execution to two coupled tasks: \emph{input scheduling} (choosing which tuples enter each batch at minimal header and resubmission cost) and \emph{output digestion} (reusing returned local permutations to certify prefixes and stop early on demand).

Existing listwise \textsf{ORDER BY} schedules follow three routes, yet none satisfies all requirements of a database sort operator. First, \emph{pivot partitioning} in ListK~\cite{DBLP:journals/corr/abs-2603-17223} expends window slots on pivots and coarse bucketing rather than intra-bucket order, depends on pivot sampling, outputs an unordered set in top-$K$ selection, and models only macro-pipeline recall--latency rather than intra-operator token plans. Second, \emph{external merge} in OrderByLLM~\cite{DBLP:journals/corr/abs-2509-00303} preserves local permutations as sorted runs, but hardcodes two-way half-window materialization ($k{=}2, h{=}\lfloor m/2\rfloor$) with post-pass $K$-truncation, overlooking multi-way planning and demand-driven streaming. Third, \emph{preference-graph completion} in TBS~\cite{akinola2026tbs} stores revealed pairs in a DAG to skip transitively implied batches. However, maintaining a global DAG does not yield an optimal listwise subset-cover or token-minimal plan: (i) its fixed-$m$ grid and myopic greedy packing ($\max_{c\notin B}|U[c]\cap B|$) share the $\Omega((n/m)\log_m n)$ call bound with multi-way merge (outperforming only $k{=}2$ merge in prior study), miss cross-chain transitive dividends during packing, issue underfilled residual batches that pay full header cost $H$, and resubmit all known items whenever a batch retains even one unresolved pair; (ii) its blocking pipeline requires global topological sort at termination, unable to certify streaming \LIMIT~$K$ prefixes even under coverage thresholds; (iii) it lacks a label-free token planner, tuning $m$ via ground-truth sweeps; and (iv) under noisy judgments, cycle rejection can stall on conflicts and propagate early wrong edges along transitive chains.

We therefore redesign multi-way external merge across both intra-operator state mechanics and pre-execution planning. Because sorted runs encode transitive chains implicitly, a merge group need only inspect enough leading rows per run to certify the next output slice while isolating unaffected runs. \emph{Sliding-Window Merge} (SW) maintains a fixed-quota frontier across $k$ runs, emitting certified safe prefixes and retaining unresolved survivors. \emph{Segmented Loser-Tree Merge} (LT($h$)) stores losing $h$-row segments and refill pointers inside a binary tournament tree, refreshing only one root-to-leaf path of $\lceil\log_2 k\rceil + 1$ segments per call while leaving off-path branches untouched. Both operators unify pass-materialized \emph{Full} execution with a root-driven pull \emph{producer tree} for \LIMIT~$K$. As Figure~\ref{fig:framework} illustrates on 54 Movie Q9 reviews~\cite{lao2026sembench}, pricing both headers ($H$) and item appearances ($L$) leads our planner to pick $\mathrm{SW}(3,2)$ for Full (5,350 merge input tokens) but single-pass $\mathrm{LT}(9,1)$ for \LIMIT~3 (589 tokens).

\stitle{Contributions.} Our key contributions are summarized as follows:
\begin{itemize}[leftmargin=*,nosep,topsep=2pt]
  \item \textbf{Multi-Way Listwise Merge and Streaming \LIMIT~$K$ (Section~\ref{sec:framework}).} We design SW with safe-prefix certification and LT($h$) with pointer-guided path refresh, unifying materialized Full passes with root-driven demand-pull \LIMIT~$K$ execution and proving sortedness, prefix equivalence, and run-order guarantees.
  \item \textbf{Label-Free Runtime Input-Token Planner (Section~\ref{sec:cost}).} We model merge input tokens as $\widehat{C}(\theta)=H\widehat{\mathbb{E}}[A]+L\widehat{\mathbb{E}}[S]$, profiling $(H,L)$ locally without labels or LLM calls and replaying cached finite schedules to jointly select $(f,k,h)$ across demands and prefix caching.
  \item \textbf{Cross-Paradigm Evaluation and Diagnosis (Section~\ref{sec:exp}).} Across six benchmarks and two model scales, we evaluate quality, tokens, and latency alongside $(k,h)$ cost surfaces, planner regret, scalability to $10\text{K}$ rows, prefix-cache pricing, and limit sweeps.
\end{itemize}
